\section{Conclusion}
\label{sec:conclusion}

We began by observing that a practitioner applying CCNet's standard perplexity filter to RedPajama-V2 would retain 86\% of Spanish documents but only 16\% of English documents --- not from any quality difference, but from a calibration artifact introduced by applying a global threshold to locally-calibrated perplexity scores. Having characterized this phenomenon precisely, we can now say what kind of bias it is: a large, structural, language-family divide (Cram{\'e}r's $V = 0.40$--$0.57$) that is consistent across the full practical range of filtering aggressiveness, maps exactly onto the Romance--Germanic divide, and is 25--40\% larger than prior estimates suggested.

Our contributions are: (1) the first calibrated $V$ baseline for global \texttt{ccnet\_perplexity} thresholding on RedPajama-V2, (2) quantification of the language-family structure of the disparity, and (3) recalibration of prior effect size estimates that will affect how researchers design correction studies.

The most urgent next step is testing whether per-language $k$-th percentile calibration reduces $\Delta V \geq 0.10$ for $\geq 3$ of 5 $k$ values --- the primary correction hypothesis. If confirmed, per-language calibration should become the default, not the exception, in multilingual dataset curation. The challenge of calibrating multilingual quality filters equitably has been recognized qualitatively for years. This work provides the quantitative foundation needed to evaluate progress: a precise, reproducible $V$ measurement that practitioners and researchers can use as the standard for multilingual dataset curation going forward.
